%% UML — Diagrama de clases (plantilla)
%% Compilar: tectonic -X compile uml-clases.tex
\documentclass[11pt,a4paper]{article}
\usepackage{../../docs/latex/legasoft}
\usepackage{../../docs/latex/legasoft-diagramas}

\lgtitulo{UML: diagrama de clases}
\lgtituloCorto{UML — Clases}
\lgcodigo{LGS-MOD-UML-2}
\lgversion{0.1}
\lgestado{Plantilla}
\lgfecha{\campo{fecha}}
\lgautor{\lgArquitecto\ (\lgArquitectoRol)}

\begin{document}

\begin{center}
{\sffamily\bfseries\LARGE\color{lgPrimario} UML: diagrama de clases\par}
\vspace{2pt}
{\sffamily\normalsize\color{lgGris} \campo{nombre del sistema} · Legasoft\par}
\end{center}
\vspace{6pt}
{\color{lgBorde}\hrule height 0.8pt}
\vspace{10pt}

\begin{porcompletar}[Cómo usar esta plantilla]
Las clases mostradas son un ejemplo de la notación, no un modelo del dominio del
cliente. Sustitúyalas por las entidades reales una vez definido el alcance en
\texttt{LGS-REQ-001}. Visibilidad: \texttt{+} pública, \texttt{-} privada,
\texttt{\#} protegida.
\end{porcompletar}

% =============================================================================
\section*{Diagrama}

\begin{figure}[H]
\centering
\ajustaancho{%
\begin{tikzpicture}[uml, node distance=1.6cm and 1.8cm]

  % --- Clases ---------------------------------------------------------------
  \node[umlclase] (usuario) {%
    \campo{Usuario}
    \nodepart{two}
      - id: int \\
      - nombre: string \\
      - correo: string \\
      - rol: Rol
    \nodepart{three}
      + autenticar(): bool \\
      + tienePermiso(p): bool};

  \node[umlclase, right=2.6cm of usuario] (registro) {%
    \campo{Registro}
    \nodepart{two}
      - id: int \\
      - fecha: date \\
      - estado: Estado \\
      - detalle: string
    \nodepart{three}
      + validar(): bool \\
      + cambiarEstado(e) \\
      + total(): decimal};

  \node[umlclase, below=2.2cm of registro] (linea) {%
    \campo{LineaRegistro}
    \nodepart{two}
      - id: int \\
      - cantidad: int \\
      - precio: decimal
    \nodepart{three}
      + subtotal(): decimal};

  \node[umlclase, below=2.2cm of usuario] (admin) {%
    \campo{Administrador}
    \nodepart{two}
      - alcance: string
    \nodepart{three}
      + crearUsuario(u) \\
      + desactivar(u)};

  % --- Relaciones -----------------------------------------------------------
  % Asociación Usuario 1 --- * Registro
  \draw[umlasoc] (usuario.east) -- (registro.west)
    node[umlmult, above, pos=0.08] {1}
    node[umlmult, above, pos=0.92] {0..*}
    node[font=\sffamily\scriptsize, text=lgGris, below, pos=0.5] {registra};

  % Composición Registro <>--- * LineaRegistro
  \draw[umlcomposicion] (registro.south) -- (linea.north)
    node[umlmult, right, pos=0.12] {1}
    node[umlmult, right, pos=0.88] {1..*};

  % Herencia Administrador ---|> Usuario
  \draw[umlherencia] (admin.north) -- (usuario.south);

\end{tikzpicture}}
\caption{Clases del dominio, sus atributos, operaciones y relaciones.}
\end{figure}

% =============================================================================
\Needspace*{20\baselineskip}
\section*{Notación de relaciones}

\begin{center}
\begin{tabularx}{\linewidth}{@{}L{3.4cm} L{3.0cm} Y@{}}
\toprule
\sffamily\bfseries\color{lgPrimario}Relación &
\sffamily\bfseries\color{lgPrimario}Símbolo &
\sffamily\bfseries\color{lgPrimario}Significado \\
\midrule
Herencia    & Triángulo hueco   & La subclase especializa a la superclase. \\
Composición & Rombo relleno     & La parte no existe sin el todo; se elimina con él. \\
Agregación  & Rombo hueco       & La parte puede existir independientemente del todo. \\
Asociación  & Línea simple      & Vínculo estructural, con multiplicidad en cada extremo. \\
Dependencia & Línea discontinua & Uso puntual: un cambio en el destino puede afectar al origen. \\
\bottomrule
\end{tabularx}
\captionof{table}{Relaciones UML empleadas en los modelos de Legasoft.}
\end{center}

\begin{nota}[Relación con el modelo de datos]
El diagrama de clases describe el dominio en términos de objetos. El modelo de
datos relacional que lo persiste se documenta por separado en
\texttt{LGS-ARQ-001}; ambos deben mantenerse coherentes, pero no son el mismo
artefacto.
\end{nota}

\end{document}
